\section{Model description}
\label{sec:model}

The model is a compact one-dimensional convolutional network that reads the
raw AIM-2 sensor signals and outputs the probability that an 8\,s window
contains eating (Figure~\ref{fig:architecture}). It has four parts. A
convolutional feature extractor learns its own filters from the time series
instead of receiving hand-crafted features or a spectrum. A squeeze-and-excitation
gate reweights the learned feature channels in each window. An attention
pooling layer weights the time steps that carry the evidence. Finally, the
output has two heads: one classifies eating, and one counts chews and is used
only during training. The network has 47,282 parameters at inference
(Table~\ref{tab:layers}). Table~\ref{tab:model} lists every design decision
with its reason and the literature it rests on.

\begin{figure}[H]
\centering
% Model architecture, top to bottom. Grayscale only; font inherited from the
% document (Arial 11). Drawn at final size - never scale it, or the text stops
% being 11 pt. Used by literature_review.tex and figure_architecture.tex.
\begin{tikzpicture}[
    box/.style   = {draw=black!70, fill=black!5, rounded corners=1.5pt,
                    align=center, inner sep=4pt, text width=78mm, execute at begin node={\hyphenpenalty=10000},
                    minimum height=12mm},
    key/.style   = {box, fill=black!17},
    head/.style  = {box, text width=37mm},
    flow/.style  = {-{Stealth[length=2mm]}, draw=black!75, line width=0.6pt},
    dims/.style = {text=black!60, font=\itshape, anchor=west},
    stage/.style = {text=black!60, anchor=east, align=right, execute at begin node={\hyphenpenalty=10000}},
    node distance = 4mm,
  ]
  \node[box] (in)
      {\textbf{Input}: 8 s window at 128 Hz, accelerometer X, Y, Z and optical};
  \node[box, below=of in] (norm)
      {\textbf{Normalise}: per window and channel, mean removed, z-score over valid samples, gaps → 0};
  \node[box, below=of norm] (c1)
      {\textbf{Conv block 1}: 32 filters, kernel 9, stride 2, batch norm, ReLU, max-pool 2};
  \node[box, below=of c1] (c2)
      {\textbf{Conv block 2}: 64 filters, kernel 9, stride 2, batch norm, ReLU, max-pool 2};
  \node[box, below=of c2] (c3)
      {\textbf{Conv block 3}: 64 filters, kernel 5, batch norm, ReLU};
  \node[key, below=of c3] (se)
      {\textbf{Channel attention} (squeeze-excite): global average → 16 → 64, sigmoid gate rescales each feature};
  \node[key, below=of se] (ap)
      {\textbf{Attention pooling over time}: one score per step, softmax over 64 steps, weighted sum};
  \node[box, below=of ap] (fc)
      {\textbf{Dense}: 64 units, ReLU, dropout 0.3};
  \node[head, below=7mm of fc.south west, anchor=north west] (food)
      {\textbf{Eating head}\\sigmoid → p(eating)};
  \node[head, dashed, fill=white, below=7mm of fc.south east, anchor=north east] (chew)
      {\textbf{Chew-count head}\\32 → 1, softplus};

  \foreach \a/\b in {in/norm, norm/c1, c1/c2, c2/c3, c3/se, se/ap, ap/fc}
    \draw[flow] (\a) -- (\b);
  \draw[flow] (fc.south -| food.north) -- (food.north);
  \draw[flow, dashed] (fc.south -| chew.north) -- (chew.north);

  % output shapes, right of each block
  \foreach \n/\s in {in/1024 × 4, norm/1024 × 4, c1/256 × 32, c2/64 × 64,
                     c3/64 × 64, se/64 × 64, ap/64, fc/64}
    \node[dims] at ([xshift=3mm]\n.east) {\s};
  \node[dims, align=left] at ([xshift=3mm]chew.east) {training only,\\loss weight 0.2};

  % stages, left of the column
  \node[stage] at ([xshift=-3mm]$(in.west)!0.5!(norm.west)$) {Signal};
  \node[stage] at ([xshift=-3mm]c2.west) {Feature\\extractor};
  \node[stage] at ([xshift=-3mm]$(se.west)!0.5!(ap.west)$) {Attention};
  \node[stage] at ([xshift=-3mm]food.west) {Output,\\threshold fit on\\validation\\participants};
\end{tikzpicture}

\caption{Model architecture. Shaded blocks are the two attention mechanisms.
The dashed chew-count head is trained jointly and discarded at inference.
Output shapes are time steps × features.}
\label{fig:architecture}
\end{figure}

% Layer-by-layer dimensions and parameter counts (inference model plus the
% training-only chew head). Counts computed from the layer definitions in
% scripts/train_food_intake.build_attention_cnn.
\begin{table}[H]
\centering
\caption{Layer dimensions. An 8\,s window at 128\,Hz is 1024 samples; each of
the 64 output time steps covers about 0.9\,s of signal at 125\,ms spacing.}
\label{tab:layers}
\begin{tabular}{@{}llrr@{}}
\toprule
\rowcolor{black!8}
\textbf{Layer} & \textbf{Configuration} & \textbf{Output} & \textbf{Parameters} \\
\midrule
Input            & Acc X, Y, Z, Optical            & 1024 × 4  & -- \\
Conv block 1     & 32 × \textit{k}9, stride 2, BN, max-pool 2 & 256 × 32  & 1,312 \\
Conv block 2     & 64 × \textit{k}9, stride 2, BN, max-pool 2 & 64 × 64   & 18,752 \\
Conv block 3     & 64 × \textit{k}5, BN                       & 64 × 64   & 20,800 \\
Channel attention & 64 → 16 → 64, sigmoid          & 64 × 64   & 2,128 \\
Attention pooling & 1 × 1 conv, softmax over time  & 64        & 65 \\
Dense            & 64, ReLU, dropout 0.3           & 64        & 4,160 \\
Eating head      & 1, sigmoid                      & 1         & 65 \\
\midrule
\multicolumn{3}{@{}l}{\textit{Total at inference}}               & \textit{47,282} \\
Chew-count head  & 32 → 1, softplus (training only) & 1        & 2,113 \\
\bottomrule
\end{tabular}
\end{table}


% Complete model description, one row per design decision.
% Style follows the 9/22 lab-meeting table: #, component, choice, why - plus
% the literature each decision rests on.
\begin{xltabular}{\textwidth}{@{}>{\raggedright\arraybackslash}p{5mm}
                               >{\raggedright\arraybackslash}p{24mm}
                               >{\raggedright\arraybackslash}X
                               >{\raggedright\arraybackslash}X
                               >{\raggedright\arraybackslash}p{24mm}@{}}
\caption{Complete description of the food-intake detection model. Every design
decision, the reason for it, and the literature it is grounded in.}
\label{tab:model}\\
\toprule
\rowcolor{black!8}
\textbf{\#} & \textbf{Component} & \textbf{Design} & \textbf{Why} & \textbf{Literature} \\
\midrule
\endfirsthead
\toprule
\rowcolor{black!8}
\textbf{\#} & \textbf{Component} & \textbf{Design} & \textbf{Why} & \textbf{Literature} \\
\midrule
\endhead
\midrule
\multicolumn{5}{r@{}}{\textit{continued on next page}} \\
\endfoot
\bottomrule
\endlastfoot

\multicolumn{5}{@{}l}{\textit{Data}} \\
1 & Sensor & AIM-2 eyeglass sensor: 3-axis accelerometer and optical temporalis-muscle sensor, 128\,Hz, stored as 8\,s packets in the ETS database & Chewing moves the temporalis muscle and the head; both channels carry it, with different artefacts. The original AIM-2 used a flex sensor where this one has the optical sensor & \citet{doulah2021aim2, ghosh2022comparative} \\
2 & Labels & Eating = bout OR bite annotation at 10\,Hz. Window is eating if more than half its annotated samples are eating. Unrecorded time (−1) is dropped, never treated as non-eating & Treating unannotated time as negative would bury the real labels and invent easy negatives & (this work) \\
3 & Alignment & No time shift. Sensor/annotation lag measured by cross-correlating the optical chewing envelope with the labels: median +0.10\,s, all within 0.2\,s on the 20 sessions measured first (to be repeated on all 61) & A correction must be measured, not assumed; a residual lag of one packet would corrupt every label edge & (this work) \\
4 & Window & 8\,s (1024 samples), non-overlapping at evaluation & Chewing lies at 1--3\,Hz; a 5\,s window was chosen for chewing detection so that it holds at least five chews, and AIM-2 used 10\,s epochs. An 8\,s window holds at least eight chews and matches the packet size. Lengths of 2--16\,s are also evaluated & \citet{diou2022intake, doulah2021aim2, po2011chewing} \\

\multicolumn{5}{@{}l}{\textit{Input}} \\
5 & Representation & Normalised time series, 1024 samples × 4 channels; no spectrum and no hand-crafted features & In the reviewed studies, features learned from the signal outperformed hand-engineered ones. None compares time-series and spectral input on the same network; that comparison is made here (61 sessions, 5-fold): F1 0.858 against 0.796, and 32 against 81 false alarms per hour of non-eating & \citet{ghosh2022comparative, kyritsis2021endtoend, diou2022intake}; this work \\
6 & Normalisation & Per window, per channel: mean removed, z-scored over valid samples only; missing samples set to 0 after scaling & Optical baseline varies by participant and fit on the face; DC level is not information about chewing. A literal −1 against counts near 3000 injects a step at every gap. z-normalisation is likewise the only preprocessing in the standard convolutional time-series baselines & \citet{wang2017strong} \\

\multicolumn{5}{@{}l}{\textit{Network}} \\
7 & Convolutional trunk & Three 1D conv blocks: 32@\textit{k}9/s2, 64@\textit{k}9/s2, 64@\textit{k}5, each with batch norm and ReLU, max-pool 2 after the first two. Output 64 time steps × 64 features & Learns its own filters from the signal. For chewing detection, a CNN reached F1 0.908 against 0.748 for an SVM on hand-crafted features. Each output step covers about 0.9\,s of signal (one chew cycle) at 125\,ms spacing & \citet{wang2017strong, ghosh2022comparative, diou2022intake, hannun2019cardiologist} \\
8 & Channel attention & Squeeze-and-excitation: global average → dense 16 (ReLU) → dense 64 (sigmoid) → rescale each feature channel & Walking at 1--3\,Hz resembles chewing and produces false positives in chewing detectors; the accelerometer is what separates the two. The gate reweights optical against motion channels in every window at negligible cost & \citet{hu2018squeeze, diou2022intake} \\
9 & Temporal pooling & Attention pooling: one score per time step (1 × 1 conv), softmax over the 64 steps, weighted sum & Global averaging dilutes a 1--2\,s bite across 8\,s. Attention weights the informative steps, as in multiple-instance learning & \citet{ilse2018attention} \\
10 & Classifier & Dense 64 (ReLU), dropout 0.3, sigmoid output & Small head; about 47\,k parameters in total at inference & -- \\
11 & Auxiliary head & Chew-count regression (dense 32 → 1, softplus), target scaled by its 99th percentile. Used in training only & Deep models require large volumes of labelled data, and one label bit per window is thin supervision. The chew-count task is specific to this work; the principle is multi-task learning & \citet{caruana1997multitask, diou2022intake} \\

\multicolumn{5}{@{}l}{\textit{Training}} \\
12 & Loss & Binary cross-entropy (weight 1.0) + chew MSE (weight 0.2); class-balanced sample weights computed per fold & Auxiliary loss is a regulariser, not the objective. Eating fraction ranges 0--86\% by session, so the rare class changes with the fold & \citet{caruana1997multitask} \\
13 & Augmentation & 50\% overlapping windows (4\,s hop) at training time only & Roughly doubles the training windows at no collection cost; evaluation stays on the non-overlapping grid so no second is scored twice & \citet{um2017augmentation} \\
14 & Optimisation & Adam, learning rate 0.0003, batch 64, up to 80 epochs, early stopping (patience 12) and learning-rate halving on validation ROC AUC & At 0.001 the validation optimum fell at the first epoch in most folds; the lower rate moves it later. Validation participants are held out from training, never random windows & -- \\

\multicolumn{5}{@{}l}{\textit{Evaluation}} \\
15 & Validation & Leave-one-subject-out: 43 folds, each holding out one participant with all of their sessions. Inner validation on further held-out participants & Windows of one meal are near-duplicates; any split below participant level leaks and overstates performance. Leave-one-subject-out is the protocol of the AIM-2 and smartwatch intake studies & \citet{doulah2021aim2, kyritsis2021endtoend, diou2022intake, saeb2017need} \\
16 & Threshold & Chosen on the inner validation participants, never on the test participant: at the centre of the F1 plateau by default, or as the highest recall within a false-alarm budget per hour & The F1 curve is flat near its peak, so its exact argmax on few participants is noise. A detector can be tuned towards recall or precision according to the use case; in free living, false alarms carry the higher cost & \citet{diou2022intake} \\
17 & Post-processing & None. Median smoothing over 3 windows was tested and rejected (F1 −0.079) & Meals are punctuated by pauses; at 8\,s the label genuinely switches, and smoothing erases short real bouts & (this work) \\
18 & Metrics & F1 and balanced accuracy lead; accuracy, precision, recall and false alarms per hour of non-eating reported alongside, pooled over all held-out windows and per participant & Class balance differs by participant; accuracy alone rewards predicting the majority class. False alarms per hour expresses the free-living failure mode directly & \citet{bell2020scoping} \\
\end{xltabular}


\subsection{Dataset and current performance}

The dataset is 61 annotated eating sessions (breakfast, lunch, dinner and
snack) from 43 participants, 2 to 55 minutes each: 24.2 hours of annotated
recording, of which 9.6 hours is eating. The eating fraction ranges from 0\%
to 86\% by session. Validation is leave-one-subject-out: 43 folds, each
holding out one participant with all of their sessions.

In a preliminary participant-level 5-fold cross-validation, pooled over all
10,891 held-out non-overlapping windows, the model reaches F1 0.858,
accuracy 0.891, balanced accuracy 0.881, precision 0.884 and recall 0.834,
with 32 false alarms per hour of non-eating. Per-participant F1 has mean
0.835 (standard deviation 0.136, median 0.878); the spread between people is
the honest measure of how the model generalises to someone new, and it is
larger than any difference between design choices seen so far.
Leave-one-subject-out results will replace these figures.
